\section{Hardness under degree, size, and data restrictions}\label{s3:hardness}

The main results of this section show that three plausible sources of
tractability are insufficient: degree two at every layer, two pools and
two outputs, and a single scalar mixing variable with bounded local data.
The restrictions must be read together. A network with few mixing arcs
can still have many external sources linked by bypasses, and small
physical coefficients can encode precise relationships through a large
network.

The degree results sharpen the earlier classification of
\citet[Sections~4--6]{haugland2016-the-computational-complexity-of-the}.
That classification gives scalar-quality hardness and separate hardness
results for bounded in-degrees or bounded out-degrees, without combining
all these restrictions. \citet[author manuscript, Section~4, p.~8]{s4:boland2017}
explicitly ask whether the scalar-quality case becomes polynomial when
all in-degrees, or all out-degrees, are at most two.
\Cref{s3:all-two-thm} answers both questions negatively in a common
restriction class: every one of the four layer degrees is exactly two,
capacities are one, only endpoint upper qualities are used, and profit
has no PTAS unless $\PP=\NP$. The orientation construction first develops
the two pure modes; the independent-set construction then imposes all
four degree restrictions simultaneously.

The fixed-count results address a different question.
\citet[Section~6, p.~214]{haugland2016-the-computational-complexity-of-the}
and \citet[author manuscript, Section~4, p.~8]{s4:boland2017} ask which
one-pool tractability results extend to two pools.
\Cref{s3:two-pools} rules out the extension based only on bounding the
number of outputs, even on an undirected tree; its attribute and input
counts grow. Allowing bypasses changes the one-pool boundary as well,
as explicitly noted by
\citet[Section~4.5, p.~607]{s3:haugland-hendrix2016}.
\Cref{s3:constant-thm} gives strong NP-completeness with one upper scalar
quality, two actual feeds, two actual outlets, and fixed finite physical
data. The many bypass sources encode a linear system; they are not
additional mixing inputs.

Broad one-pool capacity hardness was already asserted in
\citet[Remark~4.6]{s3:baltean2018}, whose polynomial result assumes
unbounded feed availability and pool capacity and fixed product demands.
Our contribution is the explicit reduction with the simultaneous
restrictions just stated. To the best of our knowledge, neither that
finite-data restriction nor the simultaneous scalar all-degree-two
restriction was established in the cited pooling complexity literature.
We claim no priority for general one-pool hardness or for the
positive-product source problem used in the reductions.

The proof sequence separates the mechanisms behind these results.
After the degree constructions, a positive-product objective is realized
by two outputs. Copy networks replace its attribute rows by physical
bypass constraints. An explicit penalty then removes contracts while
preserving optimal value. Finally, binary circuits and a saturation
threshold give the stronger finite-data decision result, and an exact
contract version leaves only five exceptional external nodes.
All networks are acyclic and have no pool-to-pool arcs. Degrees count
physical arcs, including bypasses where present.

We use profit, the negative of cost, when discussing approximation. A ratio
for a possibly negative minimum cost would not express the guarantees below.
The NP membership statements for fixed mixing parameters invoke
\cref{s2:pooling-np}; they do not assume rational optimal flows. For
input qualities in $[0,1]$ and only upper output quality bounds in
$\{0,1\}$, the disjunctive certificate \eqref{s1:endpoint-or} supplies
NP membership for threshold decision even with many pools.

\subsection{Orientation and the first degree restrictions}\label{s3:orientation}

A capacitated orientation of an undirected graph assigns each edge to one
of its endpoints. The load at a vertex is the sum of the weights assigned
to it. We refine the satisfiability construction of
\citet[Section~5]{s3:asahiro2011} by using private weight-two triangles.
The proof below includes the occurrence restriction and the bipartite
subdivision, which together give the degree-three refinement.

\begin{lemma}\label{s3:co}
Deciding whether every vertex load can be at most its capacity is
NP-complete on simple bipartite graphs of maximum degree three, with all
weights and capacities in $\{1,2\}$. All vertices of degree three can be
required to belong to one side of the bipartition.
\end{lemma}
\begin{proof}
Start with a 3-SAT instance whose clauses each use three distinct
variables, a standard NP-complete restriction. For a variable occurring $r\ge2$ times,
replace its occurrences by distinct variables $x_1,\ldots,x_r$ and add
$x_i\vee\neg x_{i+1}$ cyclically. These clauses force all copies equal:
if one is false, its successor must be false, and the cycle propagates
this value. A variable occurring once needs no cycle. Each resulting
literal occurs at most twice, since a copy occurs once in an original
clause, once positively in a cycle clause, and once negatively in a cycle
clause. The clauses have two or three literals on distinct variables.
Thus satisfiability remains NP-hard under these restrictions.

Make two vertices for each variable's literals and join them by a
weight-two edge. Make a vertex for each clause and join it by weight-one
edges to its literals. A two-literal clause also has a weight-one edge to
a vertex of a new private triangle whose three edges have weight two.
Every clause vertex now has degree three. Literal vertices have degree at
most three and triangle vertices have degree at most three. Set every
capacity to two. The graph is simple.

Given a satisfying assignment, assign each variable edge to its false
literal. Orient each triangle cyclically, so every triangle vertex has
load two. For each clause choose one true literal, assign its clause edge
to that literal, and assign the other two clause edges to the clause.
A true literal receives at most two unit edges. Thus every load is at
most two. Conversely, the weight-two edges in a triangle must enter
three different vertices, filling all three capacities. Its clause edge
must therefore enter the clause. The endpoint receiving a variable edge
is called false; it too cannot receive a clause edge. Each clause must
assign at least one of its three unit edges away from itself, necessarily
to a true literal. This gives a satisfying assignment.

Finally subdivide every edge $e=uv$ once, giving both half-edges weight
$w_e$ and the new vertex $m_e$ capacity $w_e$. An edge assigned to $v$
extends by assigning $um_e$ to $m_e$ and $m_ev$ to $v$. Conversely, both
half-edges cannot enter $m_e$. If one does, assign the original edge to
the other original endpoint; if neither does, choose either endpoint.
No original load increases. The subdivision is bipartite, keeps original
degrees, and gives all new vertices degree two. Membership in NP follows
by listing an orientation.
\end{proof}

\begin{theorem}\label{s3:local-degrees}
One-attribute pooling threshold decision without bypasses is strongly
NP-complete under either of these restrictions: every input has out-degree
one and every output has in-degree at most three; or every output has
in-degree one and every
input has out-degree at most three. In both cases every pool has exactly
two inputs and two outputs, capacities lie in $\{1,2\}$, input qualities
and output upper bounds lie in $\{0,1\}$, and costs lie in
$\{-2,-1,0,1\}$. Only upper output quality specifications are imposed,
and there are no positive lower flow bounds.
\end{theorem}
\begin{proof}
Let $(U\cup W,E)$ be an instance of \cref{s3:co}, with capacities $T_v$
and weights $w_e$. In the first construction, edge $e=uw$ has a private
quality-zero input $a_e$, a private quality-one input $d_e$, and a pool
$\ell_e$, all of capacity $w_e$. Vertex $v$ is an output of capacity
$T_v$; its upper quality is zero for $v\in U$ and one for $v\in W$.
The four arcs at $\ell_e$, in order from the clean input, from the dirty
input, to $u$, and to $w$, have costs $1,0,-2,-1$.

Write these flows as $a,d,s,b$, with throughput
$t=a+d=s+b\le w_e$. Their cost is
\begin{equation}
 a-2s-b=a-s-t.\label{s3:edge-cost}
\end{equation}
If $s>0$, the strict output's upper bound gives zero pool quality:
all terms in its quality-mass inequality are nonnegative. Hence $d=0$
and $a=t\ge s$. If $s=0$, $a\ge s$ is immediate. The edge cost is
therefore at least $-w_e$. Equality in the total lower bound
$-\sum_e w_e$ forces $t=w_e$ and $a=s$ at every pool. If $s>0$, the
pool uses $a=s=w_e,d=b=0$; otherwise it uses $d=b=w_e,a=s=0$.
Assign $e$ to the corresponding output endpoint. Its load is exactly
that output's inflow and respects $T_v$. Conversely a feasible orientation
routes $w_e$ in the corresponding pure mode and attains the bound.

For the second construction, each graph vertex is instead an input of
capacity $T_v$, quality zero on $U$ and one on $W$. Every edge has a
capacity-$w_e$ pool and two private capacity-$w_e$ outputs, with bounds
zero and one. Use the same four arc costs. Formula \eqref{s3:edge-cost}
and its equality conditions are unchanged. Now the load at an endpoint
is the total flow taken from its input. The same orientation equivalence
proves hardness. The graph degrees give the asserted physical degrees,
and all data are bounded constants apart from the linear-size threshold.
Membership in NP follows from the endpoint disjunction
\eqref{s1:endpoint-or}.
\end{proof}

These proofs also identify two distinct extensions. First, the pool
capacity can be omitted. In the private-input construction, $s>0$
forces $d=0$, so cost $=-s\ge-w_e$; when $s=0$, cost $=-d\ge-w_e$.
Equality in the first case sends $w_e$ to the strict endpoint; in the
second it takes $w_e$ from the dirty input and sends at least $w_e$ to
the lax endpoint. Thus assigning the edge accordingly never exceeds
its output's load. Conversely the earlier orientation flows remain
feasible. In the private-output construction, the same cost expressions
apply; $s\le w_e$ and, when $s=0$, $d\le b\le w_e$. At equality the
selected input supplies at least $w_e$. This again gives the orientation.
Additional cost-neutral cross-routing may now occur, so equality need
not make every physical flow integral.

Second, apply either construction to two vertices joined by parallel
edges of positive integer weights $s_1,\ldots,s_n$, with both capacities
$B$ and $\sum s_i=2B$. A feasible orientation is precisely a partition
of the weights into equal sums. This proves ordinary NP-hardness with
two outputs and private inputs, or two inputs and private outputs,
respectively; it does not establish strong hardness for these two-node
specializations. Parallel edges in the source create distinct pools,
not parallel physical arcs.

At the other end of the degree boundary,
\citet[Proposition~3]{haugland2016-the-computational-complexity-of-the}
shows that if every pool has in-degree one or out-degree one, a network
without bypasses has an LP formulation.
An in-degree-one pool has fixed quality. For any remaining pool, which
has one outlet, substitute its incoming quality mass directly in the
unique output's quality inequality. Keep its linear throughput capacity
and mass balance. A feasible LP point recovers the pool quality by the
usual mixture at positive throughput and arbitrarily at zero throughput.
This works coordinate by coordinate. Thus $(2,2)$ is the smallest common
pool-degree pattern that permits the hardness just proved.

\subsection{All four degrees two: exact modes and approximation}\label{s3:all-two}

We use the established constant-gap hardness of maximum independent set
on simple cubic graphs supplied with a proper three-edge-coloring.
The needed source is the explicit construction in
\citet[author manuscript, Section~5(A), pp.~25--26]{s3:chlebik2006}:
it produces edge-three-colored three-regular graphs and preserves the
independent-set objective and gap. The coloring is part of the constructed
instance. No particular numerical gap constant is needed here. The coloring is
essential: cubic degree alone would not supply the three disjoint
resource matchings used below.

\begin{theorem}\label{s3:all-two-thm}
One-attribute pooling threshold decision is strongly NP-complete when all
four degree counts (input out-degree, pool in-degree, pool out-degree,
and output in-degree)
are exactly two, all vertex capacities are one, qualities and upper
specifications lie in $\{0,1\}$, and costs lie in $\{-2,-1,0,1\}$.
Only upper output quality specifications are imposed. There are no
bypasses or positive lower flow bounds. Its nonnegative
profit objective on the constructed family has no PTAS unless
$\PP=\NP$. In the constructed family, every feasible solution has an integral
pure-mode replacement of at least the same profit.
\end{theorem}
\begin{proof}
Let $G=(V,E_1\mathbin{\dot\cup}E_2\mathbin{\dot\cup}E_3)$ be the
source graph, $n=|V|$. Each color class is a perfect matching. Make a
pool $\ell_v$ for each vertex. A clean input $a_e$ for $e\in E_1$
feeds its two endpoint pools, and a dirty input $d_e$ for $e\in E_3$
feeds its two endpoint pools. A strict output $s_e$ for $e\in E_2$
receives from its two endpoint pools. Initially give each pool a private
lax output $b_v$. All capacities are one. Qualities are zero for clean
inputs and one for dirty inputs; strict and lax bounds are zero and one.
The four arcs at each pool have costs $1,0,-2,-1$ as above.

Write the flows at $v$ as $a_v,d_v,s_v,b_v$. Mass balance shows that
its profit is $s_v+d_v\ge0$. If $s_v>0$, its quality must be zero and
$d_v=0$. Keep only the clean-to-strict flow
$a'_v=s'_v=s_v$, deleting the other flows. If $s_v=0$, keep only
$d'_v=b'_v=d_v$. In either case every arc flow decreases or stays fixed,
profit is preserved, and the new qualities, respectively zero and one,
satisfy all outputs. Thus any feasible flow can be made pure, with at
most one allowed mode retained per pool.

Fix those mode choices, assigning inactive pools arbitrarily. Replace
each allowed mode by an edge between its input and output. Feasible
throughputs are precisely fractional matchings in this bipartite graph;
the unit edge bound also enforces the corresponding pool capacity.
The bipartite matching polytope is integral: orient its incidence edges
from the input side to the output side to obtain a network incidence
matrix, and add integral bounds. A maximum matching therefore has size
at least the old fractional total. Route one unit in its selected modes.
This proves the claimed replacement and is algorithmic.

Let $H$ be obtained by subdividing each edge of $E_3$ twice. For
$uv\in E_3$ write its new path $u-d_u-d_v-v$. An original vertex $v$
represents its pool's clean mode and $d_v$ its dirty mode. Clean modes
conflict along $E_1$ through a shared input and along $E_2$ through a
shared output. Dirty modes conflict on $d_ud_v$ through a shared input;
$v$ and $d_v$ conflict because they are modes of one pool. These are
all conflicts. Consequently integral mode selections are independent
sets of $H$, with profit equal to cardinality.

The following identity includes a constructive recovery:
\begin{equation}
 \alpha(H)=|E_3|+\alpha(G)=n/2+\alpha(G).\label{s3:is-identity}
\end{equation}
Indeed, extend any independent set of $G$ by one internal vertex of
each subdivided edge, avoiding its selected endpoint if present. For
the reverse direction, if an independent set of $H$ contains both
original endpoints $u,v$ of an $E_3$ edge, it contains neither internal
vertex. Delete $u$ and insert $d_u$. This preserves cardinality and
independence. Since $E_3$ is a matching, the changes do not interfere.
Afterwards the selected original vertices form an independent set of
$G$, and at most one internal vertex is present per edge. Their number
is at least the original cardinality minus $n/2$. Combining this with
integral replacement proves that the pooling optimum is exactly
$n/2+\alpha(G)$. Threshold $n/2+k$ gives strong hardness.

A cubic graph has $\alpha(G)\ge n/4$, by successively choosing a vertex
and deleting it and its at most three neighbors. Therefore
$P^*=n/2+\alpha(G)\le3\alpha(G)$. A feasible profit
$P\ge(1-\varepsilon)P^*$ yields in polynomial time an independent set
of size at least
\begin{equation}
 P-n/2\ge\alpha(G)-\varepsilon P^*
 \ge(1-3\varepsilon)\alpha(G).\label{s3:gap-transfer}
\end{equation}
A PTAS would contradict the source gap. Here and in the approximation
recovery below, we use the ordinary bit model: an approximation algorithm
returns feasible rational arc flows in binary, with output length bounded
by its polynomial running time. Cleanup uses only those flows; its exact
rational operations and the subsequent matching computations take time
polynomial in their encoding length. At positive throughput, pool quality
is recovered from the incoming quality mass; at zero throughput it may
be chosen arbitrarily.

Finally merge the two private lax outputs of the endpoints of every
$E_3$ edge into one capacity-one lax output. Now there are $n$ inputs,
$n$ pools, and $n$ outputs, all with the required degree counts exactly
two. Cleanup still only decreases resource use. In pure modes the new
constraint is $d_u+d_v\le1$, already imposed by their shared dirty input.
Thus it changes neither the integral selections nor the optimum or
recovery. Each of the two undirected layer-incidence graphs is now a
disjoint union of cycles. The endpoint disjunction
\eqref{s1:endpoint-or} gives NP membership for the full stated
zero-tolerance restriction class.
\end{proof}

Hard selection of sparse active flows was already established by
\citet[Definition~2.2 and Proposition~3.1, p.~96]{s3:haugland2019-single-flow},
who imposes restrictions on the number of positive inlet or outlet flows
at each pool and proves NP-hardness of all four single-flow variants.
In \cref{s3:all-two-thm}, integral pure-mode replacement is instead a
proved property of an ordinary continuous pooling family under the
simultaneous degree and data restrictions. The next proposition extends
that replacement argument to linked weighted costs.

\begin{proposition}[A weighted integral-mode family]\label{s3:weighted}
Allow arbitrary sharing degrees, zero flow lower bounds, and nonnegative
integer source, pool, output, and (if imposed) arc capacities. Each pool
still has one quality-zero input, one
quality-one input, one strict output with bound zero, and one lax output
with bound one. Only upper output quality specifications are imposed.
For arbitrary rational $\alpha_v,\beta_v$, give its clean
intake, dirty intake, strict outlet, and lax outlet costs
$\beta_v,0,-(\alpha_v+\beta_v),-\beta_v$. An optimal solution exists
with integral throughputs and only pure modes. With unit capacities
this optimization is exactly maximum-weight rainbow matching in a
bipartite multigraph: pool $v$ contributes a clean-input/strict-output
edge of weight $\alpha_v$ and a dirty-input/lax-output edge of weight
$\beta_v$. These two edges have disjoint endpoints and the common color
$v$; at most one edge of each color may be selected.
\end{proposition}
\begin{proof}
Mass balance makes cost $-\alpha_vs_v-\beta_vd_v$. The cleanup in the
preceding proof preserves $s_v$ and $d_v$ separately, and thus preserves
cost even for signed rewards. Once one mode per pool is fixed, maximize
weighted flow between the input and output vertices, with their integer
capacities. Bound each retained mode edge by the minimum of its pool
capacity and its two physical arc capacities, omitting any arc bound
that is not imposed. These edge bounds are integral. This is a capacitated
bipartite flow polytope with an integral optimum. It improves or preserves
the previous objective and restores a feasible integral pooling point.
The full feasible set is compact in bounded quality coordinates, so an
optimum exists to which this replacement applies. For unit capacities,
the two edges of a color have disjoint endpoints because their input and
output types differ; the only further restriction is to choose at most
one of them. This is precisely the stated rainbow matching problem.
The conclusion concerns these four linked arc costs, not arbitrary
pooling costs.
\end{proof}

\begin{proposition}[Positive strict-output concentration bounds]\label{s3:tolerance}
There is a fixed rational $\delta>0$ for which replacing the strict
bounds in \cref{s3:all-two-thm} by $\delta$ preserves absence of a profit
PTAS unless $\PP=\NP$, and exact optimization remains strongly NP-hard.
The same conclusion holds with input qualities $\{1,2\}$ and output
bounds $\{1+\delta,2\}$.
\end{proposition}
\begin{proof}
For a feasible flow with profit $P$ and any $\eta>0$, separate the pools
by whether their quality $p_v$ is at most $\eta$. For rational $\eta$,
this is the exact rational test $d_v\le\eta(a_v+d_v)$ at nonzero
throughput; empty pools may be put in either group. In the first group keep
only clean-to-strict flow $\min(a_v,s_v)$ and set quality zero. Its profit
loss is at most $2d_v=2p_vt_v\le2\eta$, since
$s_v-a_v=d_v-b_v\le d_v$. In the other group keep only dirty-to-lax
flow $\min(d_v,b_v)$ and set quality one. Its loss is at most $2s_v$,
since $d_v-b_v=s_v-a_v\le s_v$. Every arc flow weakly decreases,
so the new solution is feasible with zero strict bounds. Summing the
old strict quality inequalities gives
$\sum_v p_vs_v\le\delta\sum_v s_v\le\delta n$, whence its profit is
at least $P-2n(\eta+\delta/\eta)$. Empty pools cause no change.

Let $\gamma>0$ be a fixed independent-set approximation gap: a guarantee
$(1-\gamma)\alpha(G)$ is NP-hard. Choose a positive rational
$\eta<\gamma/32$ and put $\delta=\eta^2$. Then
$P_0^*\le P_\delta^*\le P_0^*+4\eta n$.
A PTAS for the positive concentration bound, with fixed
$\varepsilon<\gamma/6$, followed
by cleanup, integral replacement, and recovery, returns a set of size
at least
\[
 (1-\varepsilon)P_0^*-4\eta n-n/2
 \ge(1-3\varepsilon-16\eta)\alpha(G)
 >(1-\gamma)\alpha(G),
\]
a contradiction. These are fixed rational data, so the same gap reduction
also proves strong hardness of exact optimization. It does not preserve
\eqref{s3:is-identity} or assert integral optima at a positive bound.
Here the bound limits concentration exactly; the argument does not allow
arbitrary numerical residuals in the flow or quality constraints.
Adding one to all qualities and bounds leaves the flow constraints
unchanged by mass conservation and gives the final assertion.
\end{proof}

\begin{example}[A positive concentration bound can destroy integral optimality]\label{s3:tolerance-example}
Take $K_4$ on vertices $0,1,2,3$, with matchings
$E_1=\{01,23\}$, $E_2=\{02,13\}$, and $E_3=\{03,12\}$.
Use the exactly-two-degree construction with strict bound
$0<\delta\le1/2$. At pools 0 and 3 take clean flow $1-\delta$ and
dirty flow $\delta$, and send one unit to the respective strict output.
At pool 1 take one dirty unit and send it to the lax output; leave pool
2 inactive. The shared dirty input for $03$ uses $2\delta\le1$, and
the other uses one. Each clean input uses $1-\delta$, each strict output
receives one unit of quality $\delta$, and the used lax output receives
one unit of quality one. Thus this is feasible and has profit
$3+2\delta$. At zero tolerance, \eqref{s3:is-identity} gives optimum
$2+\alpha(K_4)=3$. Any integral unit-capacity flow at $\delta<1$
uses only pure pools, and a strict output receives at most one unit;
it cannot receive a dirty unit. Every such integral flow is therefore
also zero-tolerance feasible and has profit at most three. Hence the
positive-tolerance feasible point strictly exceeds every integral one.
\end{example}

\subsection{A positive-product source and two pools on a tree}\label{s3:positive-product}

The next reductions use the positive-product construction of
\citet[METR95-13, Sections~2--3, pp.~2--7]{s3:matsui1996}. We record the
arithmetic needed later and supply its separation argument directly.
This matters because the arbitrary-$p$ statement of Theorem~2.1 in the
report is stronger than the estimate needed for its hardness reduction.

\begin{lemma}[The explicit positive-product family]\label{s3:matsui}
For a zero-one matrix $M$ with $n\ge5$ columns, put $p=n^{n^4}$ and let
$\Omega(M)$ consist of $(r,s)\in[0,1]^{n+n^2}$ satisfying
\[
 Mr=\mathbf1,\quad s_{ii}=r_i,\quad
 s_{ij}\le r_i,\ s_{ij}\le r_j,\ s_{ij}\ge r_i+r_j-1\quad(i\ne j).
\]
Write
\begin{align}
 X&=\sum_i p^i r_i,&Y&=\sum_{i,j}p^{i+j}s_{ij},\notag\\
 U&=2p^{4n}-p+Y+2p^{2n}X,&
 V&=2p^{4n}-p+Y-2p^{2n}X,&K&=4p^{8n}.
 \label{s3:matsui-factors}
\end{align}
Both factors are strictly positive on $\Omega(M)$, and deciding whether
$UV\le K$ somewhere on it is NP-hard. All coefficients have polynomial
binary encoding length.
\end{lemma}
\begin{proof}
The zero-one feasibility question $Mr=\mathbf1$ is NP-complete (it is
exact cover in incidence-matrix form). Adding zero columns permits
$n\ge5$ without changing that answer. Let $N=n+n^2\le n^3$ and
$h=(n^3)^{-n^3}$. At every vertex of $\Omega(M)$, an invertible active
$N$-row system has integral coefficients of absolute value at most one.
Cramer's rule and the determinant bound $N^{N/2}\le(n^3)^{n^3}$ imply
that a nonintegral coordinate lies in $[h,1-h]$.

At a vertex where $r$ is not binary, let $k$ be the largest fractional
index. If $i>k$ or $j>k$, the McCormick constraints give
$s_{ij}=r_ir_j$, so those terms cancel in $D=Y-X^2$. The diagonal
term at $k$ is $p^{2k}r_k(1-r_k)\ge p^{2k}h/2$. Each of the other
at most $n^2$ signed terms with indices at most $k$ is bounded below
by $-p^{2k-1}$. Hence
\begin{equation}
 D\ge p^{2k-1}(ph/2-n^2)>p.\label{s3:corrected-matsui-gap}
\end{equation}
For the strict inequality, $ph=n^{n^3(n-3)}$ and $n\ge5$ imply
$ph/2>n^2+1$. If $Mr=\mathbf1$ has no binary solution, every vertex
has fractional $r$; the concavity of $Y-X^2$ and convex decomposition
into vertices extend $D>p$ to the whole nonempty polytope. If the
polytope is empty the product decision is already negative.

The exact product identity is
\begin{equation}
 UV-K=(Y-p)^2+4p^{4n}(D-p).\label{s3:product-gap}
\end{equation}
Thus the no case has $UV>K$. In the yes case a binary $r$ extends by
$s_{ij}=r_ir_j$, giving $D=0$. Since $0\le Y=X^2\le n^2p^{2n}$,
\[
 (Y-p)^2\le\max\{p^2,n^4p^{4n}\}<4p^{4n+1},
\]
where $n^4<p$. Equation \eqref{s3:product-gap} now gives $UV<K$.
Finally $U\ge2p^{4n}-p>0$ and
$V\ge2p^{4n}-p-2np^{3n}>0$, since $p^n>2n+1$.
The largest exponents in \eqref{s3:matsui-factors} require
$O(n^5\log n)$ bits. If $M$ has $m$ rows, the polytope has $N=n+n^2$ variables and
$O(m+n^2)$ rows. Thus the row count is polynomial in the source size;
no bound on $m$ relative to $n$ is needed.
\end{proof}

\begin{remark}\label{s3:matsui-scope}
Estimate \eqref{s3:corrected-matsui-gap} establishes precisely the
large-$p$ case used here. It does not use the report's printed
arbitrary-positive-integer-$p$ estimate. For example, $n=5,p=2$,
$r_i=1/2,s_{ii}=1/2,s_{ij}=0$ for $i\ne j$, and
$\epsilon=1/4$ satisfy that estimate's hypotheses but give
$Y-X^2=-279$, less than its claimed lower bound $-99/2$.
The corrected argument above preserves the positive-product hardness
source and all subsequent reductions.
\end{remark}

\begin{theorem}\label{s3:two-pools}
Pooling threshold decision is NP-complete with exactly two pools and two
outputs, arbitrary input and attribute counts, upper quality bounds only,
upper flow capacities in $\{1,2\}$, no bypasses, and no positive lower
flow bounds. Every input has out-degree one and the underlying undirected
network is a tree. One pool has one input and one output. Suppressing
that pool gives the same conclusion with one mixing pool, two outputs,
and exactly one bypass arc. These are ordinary NP-completeness claims.
\end{theorem}
\begin{proof}
Begin with a bounded positive-product instance
$\min_{w\in P_0}U(w)V(w)\le K$, $K>0$, from \cref{s3:matsui}.
If $P_0$ is empty, map it to a fixed no instance: take one variable input, the two pools,
the anchor, and the two outputs with the arc pattern below, zero rewards,
quality zero everywhere, and profit threshold one. Otherwise compute
$u_{\min}=\min_{P_0}U>0$ by rational LP and replace $U,V$ by
$2U/u_{\min},u_{\min}V/2$. Their product is unchanged and $U\ge2$.
Intersect $P_0$ with $V\le K$. This keeps all yes witnesses, because
$UV\le K$ implies $V\le K/2$. An empty intersection is again mapped
to a fixed no instance. These operations have polynomial bit complexity.

For $w\in[0,1]^s$, embed the remaining polytope in a simplex by
$z_i=w_i/s$ and $z_0=1-\sum_{i=1}^s z_i$, adding $z_i\le1/s$.
Denote its linear description by
\[
 P=\{z\ge0:\mathbf1^Tz=1,\ A_rz\le d_r\ (1\le r\le R)\}.
\]
Affine functions become linear on the simplex. Write
$a(z)=U(z)-1=\sum_i a_iz_i$ and $b(z)=K-V(z)=\sum_i b_iz_i$.
Then $a\ge1,b\ge0$ on $P$, although individual coefficients may
have either sign.

For each coordinate make an input $i$ feeding only pool $L$. Give these
inputs and arcs capacity two and pool $L$ capacity two. One anchor
input feeds only pool $H$, with all three capacities one. Pool $L$
feeds outputs 1 and 2, and $H$ feeds only output 1. All outlet arcs
and both outputs have capacity one. Reward intake $iL$ by $b_i$;
all other rewards are zero. There is one distinguished attribute and
one attribute per polytope row:
\begin{center}
\begin{tabular}{lcccc}
\toprule
Attribute & Input $i$ & Anchor & Output 1 bound & Output 2 bound\\
\midrule
Distinguished & $a_i$ & $0$ & $1$ & $\max(0,\max_i a_i)$\\
Row $r$ & $A_{ri}$ & $d_r$ & $d_r$ & $d_r$\\
\bottomrule
\end{tabular}
\end{center}

For positive throughput $T$ at $L$, its intake proportions are $z_i=x_i/T$.
Write its outlet flows $y_1,y_2$ and the anchor flow $h$; total output
flows are $D_1=y_1+h,D_2=y_2$. At output 1 the row-$r$ inequality is
$(A_rz)y_1+d_rh\le d_r(y_1+h)$; output 2 gives
$(A_rz)y_2\le d_ry_2$. Since $y_1+y_2=T>0$, at least one enforces
each row, and therefore $z\in P$. The distinguished inequality is
$a(z)y_1\le D_1$, so
\[
 b(z)T=b(z)(D_2+y_1)
 \le b(z)\bigl(D_2+D_1/a(z)\bigr)
 \le b(z)(1+1/a(z)).
\]
Conversely, for any $z\in P$, choose
\begin{equation}
 y_1=1/a(z),\quad y_2=1,\quad h=1-1/a(z),\quad
 T=1+1/a(z),\quad x_i=Tz_i.\label{s3:radial-max}
\end{equation}
Both outputs are filled, all flows respect their bounds, and every
attribute constraint holds. Inactive $L$ has zero profit. Compactness
and $a\ge1,b\ge0$ therefore give the exact identity
\begin{equation}
 P^*=\max_{z\in P}b(z)(1+1/a(z)).\label{s3:two-pool-value}
\end{equation}
For each $z\in P$,
\begin{equation}
 b(z)(1+1/a(z))\ge K
 \quad\Longleftrightarrow\quad
 (K-V)U\ge K(U-1)
 \quad\Longleftrightarrow\quad UV\le K.\label{s3:threshold-product}
\end{equation}
This proves hardness.

The topology is shown in \cref{s3:tree-figure}. The undirected edges form
the path $L$--output 1--$H$--anchor, with all variable inputs and output 2
as leaves at $L$. Suppressing $H$ creates exactly the anchor bypass.
There are three actual pool-outlet arcs in this construction and at
most four in the full two-pool, two-output class. Fixing their outlet
fractions gives bounded linear fibers even with unbounded attributes;
\cref{s2:pooling-np} also covers the suppressed version's two outlet
fractions and bypass. The bounded-flow
certificate in \cref{s2:pooling-np} gives NP membership in both versions.
\end{proof}

\begin{figure}[htbp]
\centering
\begin{tikzpicture}[>=Stealth,thick,
 every node/.style={font=\small},
 source/.style={draw,rectangle,minimum width=1.7cm,minimum height=0.65cm},
 pool/.style={draw,circle,minimum size=0.8cm},
 product/.style={draw,rectangle,minimum width=1.7cm,minimum height=0.65cm}]
 \node[source] (x0) at (0,1.25) {input $0$};
 \node (dots) at (0,0.45) {$\vdots$};
 \node[source] (xs) at (0,-0.35) {input $s$};
 \node[pool] (L) at (3,0.45) {$L$};
 \node[product] (j2) at (6,1.45) {output 2};
 \node[product] (j1) at (6,-0.65) {output 1};
 \node[source] (anchor) at (0,-2) {anchor};
 \node[pool] (H) at (3,-2) {$H$};
 \draw[->] (x0) -- node[above] {$x_0$} (L);
 \draw[->] (xs) -- node[below] {$x_s$} (L);
 \draw[->] (L) -- node[above] {$y_2\le1$} (j2);
 \draw[->] (L) -- node[above] {$y_1\le1$} (j1);
 \draw[->] (anchor) -- node[above] {$h\le1$} (H);
 \draw[->] (H) -- node[below] {$h$} (j1);
\end{tikzpicture}
\caption{The two-pool tree in \cref{s3:two-pools}. Pool $L$ mixes the
variable inputs; $H$ carries the fixed-quality anchor to output 1.
Suppressing $H$ produces one mixing pool and one bypass. The primary
output capacities, together with the distinguished quality bound at
output 1, give the radial limit $T\le1+1/a(z)$.}
\label{s3:tree-figure}
\end{figure}

Each attribute in this construction can be shifted by a common rational
constant at all inputs and outputs to make its data nonnegative, then
scaled into $[0,1]$. Conservation preserves its inequalities and the
encoding remains polynomial. Nonnegative input costs and output revenues
also suffice: choose $B\ge\max(0,\max_i b_i)$, charge $B-b_i$ at
variable inputs and $B$ at the anchor, and pay $B$ at either output.
All common $B$ terms cancel. Source and pool upper bounds are redundant
given the two output capacities and may be omitted. Alternatively, exact
unit demands at both outputs preserve \eqref{s3:two-pool-value}, since
\eqref{s3:radial-max} attains every mixture's bound with both outputs
filled. None of these modifications makes the large quality and reward
coefficients bounded integers.

\begin{proposition}[An FPTAS for the represented subfamily]\label{s3:fptas}
Given a nonempty bounded rational polytope $P$ and rational affine
functions $a\ge1,b\ge0$ on $P$, a feasible point of value at least
$(1-\varepsilon)\max_P b(1+1/a)$ can be found with at most
$\lceil1/\varepsilon\rceil+1$ LPs, for rational $0<\varepsilon<1$.
This gives an FPTAS for the pooling subfamily supplied with the
representation used in \cref{s3:two-pools}.
\end{proposition}
\begin{proof}
Set $N=\lceil1/\varepsilon\rceil$. For $g=k/N$, $0\le k\le N$,
maximize $b$ over $P$ with $a\le1/g$ if $g>0$; impose no additional
row for $g=0$. Skip infeasible LPs and evaluate the true objective at
one rational optimal vertex of every remaining LP. Retain the best.
Let $z^*$ attain the original optimum and put $t^*=1/a(z^*)\in(0,1]$.
The grid value $g=\lfloor Nt^*\rfloor/N$ keeps $z^*$ feasible and
satisfies $t^*-g<1/N$. The LP optimizer $z_g$ has $b(z_g)\ge b(z^*)$
and $1/a(z_g)\ge g$. Consequently
\[
 b(z_g)(1+1/a(z_g))\ge b(z^*)(1+g)
 \ge\OPT-b(z^*)/N\ge(1-1/N)\OPT.
\]
This also covers $\OPT=0$. All LP encodings and their rational vertex
solutions have bit length polynomial in the input and $\log N$.
Recover pooling flows by \eqref{s3:radial-max}. The number of LPs is
polynomial in $1/\varepsilon$. The algorithm neither recognizes arbitrary
pooling representations of this form nor gives an FPTAS for general
two-pool, two-output pooling.
\end{proof}

\subsection{Copying bounded signals through physical bypasses}\label{s3:copy}

We next replace the many attribute rows by ordinary sources and outputs.
A port is an input-to-output arc whose source is assigned when the gadgets
are connected. A source assigned several ports has their common quality
and precisely those outgoing arcs. Equal quality labels never identify
sources unless the construction explicitly joins their supply equations.
Every requested occurrence gets its own gadget; hence repeated terms
never create parallel arcs or reuse one physical flow twice. A \emph{signal}
is the common number encoded by these distinct flows, not an extra
physical flow that can be sent to several destinations. In every gadget,
source supplies and output demands refer to ordinary node totals from
\cref{s1:mass}.

\begin{lemma}[Copy and conversion]\label{s3:copy-lemma}
For distinct rational endpoint qualities $\alpha,\beta$, there is a
pool-free gadget whose four endpoint ports have flows
$x,x,2-x,2-x$, for an arbitrary $x\in[0,2]$, with each pair consisting
of one $\alpha$ port and one $\beta$ port. Chains of such gadgets can
copy a bounded signal to arbitrarily many ports and convert it to an
actual pool intake of any nonzero prescribed quality.
\end{lemma}
\begin{proof}
Put $\gamma=(\alpha+\beta)/2$. Make outputs $A,B$, each with exact
demand four and exact quality $\gamma$. Each receives one endpoint
port of each quality, of capacity two, and one capacity-four arc from a
common quality-$\gamma$ source of exact supply four. If endpoint flows
at $A$ are $u,v$ and middle flow is $m$, subtracting $\gamma$ times
mass balance from quality balance gives
$(\alpha-\gamma)u+(\beta-\gamma)v=0$, so $u=v=x$ and $m=4-2x$.
At $B$ the same argument gives $u=v=x'$ and $m=4-2x'$. Exact middle
supply implies $x+x'=2$. Conversely these formulas fill every constraint
for any $x\in[0,2]$.

Join the $B$-$\alpha$ port of one gadget to the $A$-$\alpha$ port of
the next with a quality-$\alpha$ source of exact supply two. It forces
$(2-x)+x'=2$, so the signals agree. Ordinary gadgets with endpoints
$(0,2)$ provide repeated positive and complementary quality-two ports.
Unused ports have distinct private sources with upper supply equal to
the port capacity and lower supply zero. End a chain with a $(0,a_i)$
gadget. Its $B$-$a_i$ port and a new pool intake share a quality-$a_i$
source of exact supply two. As the former carries $2-x_i$, the actual
intake is $x_i$. The unused $A$-$a_i$ port is supplied privately.
All reverse extensions follow from the displayed formulas.
\end{proof}

\begin{lemma}[Bounded-coefficient rows]\label{s3:row-copy}
A homogeneous system $Cx\le0$, with integer coefficients and
$x\in[0,2]^m$, has a copy-network realization of size polynomial in
$m$ and $\sum_{r,i}|C_{ri}|$. Its projection onto the original signals
is exactly that system.
\end{lemma}
\begin{proof}
For a row $c$, allocate $c_i$ quality-two positive ports $x_i$ when
$c_i>0$ and $-c_i$ complementary ports $2-x_i$ when $c_i<0$.
Assign all these ports to one quality-two source with upper supply
$B=2\sum_{c_i<0}(-c_i)$ and lower supply zero. Its total outflow is
$c^Tx+B$, so its supply bound is precisely $c^Tx\le0$. Omit zero rows.
Allocate enough copies for every occurrence and fill unused ports
privately. The copy identities prove one projection direction; filling
the chains at any feasible signal vector proves the other. Two rows
encode an equality.
\end{proof}

\begin{theorem}[One pool and fixed attribute dimension]\label{s3:base-bypass}
With bypasses, pooling threshold decision is NP-complete with one pool,
one physical quality having lower and upper output specifications, and
two pool outlets. Exact supplies and demands are allowed. Equivalently,
two attributes with upper specifications only suffice. All qualities may
be normalized into $[0,1]$. This statement is ordinary NP-completeness.
\end{theorem}
\begin{proof}
Use \cref{s3:matsui} without the LP normalization of
\cref{s3:two-pools}. Put $s=n+n^2$ and use simplex coordinates for
$w=(r,s_{ij})\in[0,1]^s$ as before. If $T=\sum_i x_i>0$, homogenizing
the simplex polytope at $z=x/T$ gives integer rows $Cx\le0$ of
polynomially bounded coefficients. For example,
\[
 sx_h\le T,\qquad sx_{ij}\ge sx_i+sx_j-T,
 \qquad s\sum_i M_{ri}x_i=T.
\]
Each row has coefficient magnitude at most $s+1$ and there are
polynomially many rows, so \cref{s3:row-copy} has polynomial size.

At the unrestricted simplex generators $w=0$ and $w=se_h$, write the
values of the factors as $U_i,V_i$. With $u_0=2p^{4n}-p$,
\begin{equation}
 u_0\le U_i\le u_0+sp^{2n}+2sp^{3n},\qquad
 V_i\le u_0+sp^{2n}<K.\label{s3:generator-bounds}
\end{equation}
The lower bound follows because all nonconstant coefficients of $U$
are nonnegative; the upper bounds follow by evaluating its $r$ or $s$
coefficient. Since $s<p^{2n}$ and $p\ge2$, the last strict inequality
holds. Thus the simplex coefficient vectors
$a_i=U_i-1>1$ and $b_i=K-V_i>0$ have polynomial bit length, although
large magnitudes.

Copy every signal $x_i$ and convert it to a quality-$a_i$ intake of
one capacity-two pool. Enforce all cone rows by \cref{s3:row-copy}.
The pool has two capacity-one outlets to capacity-one primary outputs.
The first also receives a capacity-one quality-zero anchor bypass and
has upper quality one; the second has redundant upper quality
$\max_i a_i$. Reward the actual intake $x_i$ by $b_i$. All other
arcs have zero reward.

For $T>0$ the copied rows enforce $z=x/T$ in the simplex image of
$\Omega(M)$, and the pool quality is $a(z)$. The argument of
\cref{s3:two-pools} gives profit at most $b(z)(1+1/a(z))$; conversely
\eqref{s3:radial-max} extends to the exact copy network. The zero-intake
assignment satisfies every homogeneous row and also fills the copies,
so the exact optimum, even for empty $\Omega(M)$, is
\begin{equation}
 \max\bigl(\{0\}\cup
 \{b(z)(1+1/a(z)):z\in P\}\bigr).\label{s3:copy-value}
\end{equation}
Threshold $K>0$ and \eqref{s3:threshold-product} prove hardness.
The fixed pool-quality parameter gives NP membership by
\cref{s2:pooling-np}.

Exact scalar qualities use both a lower and an upper specification.
Adding the negative of the scalar as a second coordinate expresses the
same restrictions by upper bounds only. Coordinate shifts and positive
scaling normalize them into $[0,1]$ without changing any feasible flow.
\end{proof}

For an input-cost realization of the same profit, reward the private
input serving the conversion gadget's $A$-$a_i$ port by $b_i$, and give
the actual intake zero reward. The copied port carries $x_i$, so the
objective is unchanged. A common charge and revenue at every input and
output cancel by mass conservation and make all input costs and output
revenues nonnegative. This observation will require an additional error
estimate after contracts are relaxed; equality of those two flows is
not assumed in a relaxed copy network.

\begin{lemma}[Full and half ports and bounded degrees]\label{s3:half}
The construction of \cref{s3:base-bypass} can have input out-degree at
most three, output in-degree at most three, and all flow upper bounds at
most four. Exact supplies and demands and exact scalar qualities remain
allowed. A full-port-only variant has input out-degree at most four.
\end{lemma}
\begin{proof}
First, with ordinary endpoint qualities $(0,2)$, enforce
$2v=u+w$ through an exact-supply-four source on four full ports
$v,v,2-u,2-w$. For a row from \cref{s3:row-copy}, list its $m$ port
values $f_j\in[0,2]$ and upper sum bound $B$. Pad by zeros to
$N=2^{\lceil\log_2m\rceil}$ leaves and impose pairwise averages in a
complete binary tree. A zero signal is fixed by a zero-upper-supply
source on its positive port. Induction gives root
$v_{\rm root}=\sum_jf_j/N$; a private positive root port with upper
supply $B/N$ enforces the row. When $m=1$, use the leaf itself; omit
$m=0$. Since $m\le N<2m$ and $0\le B/N\le2$, this construction has
linear size per row, upper capacities at most four, and input degree
at most four. Each gadget output still has three inputs.

The full-port construction above uses only its ordinary row ports in
the averaging trees; its conversion gadgets remain at the original
signals. For degree three, use full gadgets with endpoints zero and three,
midpoint quality $3/2$, demand four at each output, and middle supply
four. Their zero and quality-three port pairs are both $(x,2-x)$.
A half gadget has endpoints zero and three, middle quality one,
demand three at each output, and middle supply three. Give its zero
ports capacity two and its quality-three ports capacity one. At one
output the equations $u+v+m=3$ and $3v+m=3$ give $u=2v$.
Summing both outputs and the middle supply gives the port formulas
\begin{equation}
 (A_0,B_0)=(x,2-x),\qquad
 (A_3,B_3)=(x/2,1-x/2).\label{s3:half-ports}
\end{equation}
Their middle flows are $3-3x/2$ and $3x/2$, so all $x\in[0,2]$
are feasible. The zero-port formulas coincide with those of full and
conversion gadgets, so exact-two zero-quality links copy the same
signal along mixed chains.

Use an exact-supply-two quality-three source on
$v,1-u/2,1-w/2$. Its three-port equation is exactly $2v=u+w$.
If a child is $2-x$, the needed complementary half port is $x/2$.
Replace every averaging source in the same padded trees by this gate,
retaining the root bound $B/N$. Every feasible original row assignment
extends by its arithmetic averages; conversely induction recovers the
row from every feasible physical network. Allocate a distinct gadget
for each occurrence. Middle, chain, and conversion sources have degree
two, private and root sources degree one, and averaging sources degree
three. Outputs have degree three, except the two primary outputs of
degrees two and one. The pool retains its two outlets. Size and bounds
are as claimed. The dyadic root capacities have polynomial bit length;
large factor qualities and rewards still prevent a strong-hardness claim.
\end{proof}

\begin{lemma}[Closed cycles with one upper quality]\label{s3:cycles}
The degree-three realization can use one scalar quality and upper output
specifications only. It can further have input out-degree at most two
and output in-degree at most three, with flow upper bounds at most four.
Exact flow contracts remain in this lemma.
\end{lemma}
\begin{proof}
A full gadget with endpoints $0,\beta$, $\beta>0$, retains demands
four and middle supply four but only the upper quality $\beta/2$.
Writing zero and other endpoint flows as $u_j,v_j$, its inequality is
$v_j\le u_j$. For one signal, place all full gadgets, including its
conversion gadget if present, in a closed zero-port cycle: each link
has quality zero, exact supply two, and joins the second zero port to
the next gadget's first zero port. For $r$ gadgets, links supply total
zero-port flow $2r$. Output demands minus middle supplies give total
endpoint flow $4r$, so total other-port flow is $2r$. Every
$v_j\le u_j$ is therefore tight. The full formulas follow, and each
link forces the common signal to agree with its successor. This works
when different gadgets have different positive $\beta$, because their
normalized inequalities are identical. A one-gadget cycle joins two
distinct outputs and creates no parallel arc.

Put the half gadgets for a signal in a separate closed zero-port cycle.
Here the upper inequality is $2v_j\le u_j$. Links supply total zero
flow $2r$, while total endpoint flow is $3r$. Thus other-port flow is
$r$ and every inequality is tight, giving \eqref{s3:half-ports}.
Do not mix full and half gadgets in one cycle. Instead, couple their
signals by a quality-three exact-supply-two source on one extra full
positive port and two distinct complementary half ports:
\[
 x_f+(1-x_h/2)+(1-x_h/2)=2.
\]
It forces $x_f=x_h$. Every signal with a half occurrence gets these
three reserved occurrences, creating its full cycle if needed. Signals
without half occurrences need no coupling. Every common signal in
$[0,2]$ extends to both cycles and the coupling source by the formulas,
so the earlier row projection and objective identity remain exact.

The only degree-three inputs are now averaging and coupling sources.
Each has quality three, exact supply two, and port capacities
$(c_1,c_2,c_3)=(2,1,1)$. Replace it by three same-quality inputs of
exact supplies $c_h$. Input $h$ feeds its original port with flow $f_h$
and a new common collector with flow $c_h-f_h$, both arc capacities
$c_h$. Give the collector exact demand $\sum_hc_h-2=2$ and redundant
upper quality three. This is feasible exactly when $\sum_hf_h=2$,
which was the original source equation. It reduces input degree to two
and creates an output of degree three. All other inputs already have
degree at most two. The transformation has linear size and preserves
all original signal and intake flows. All quality values are nonnegative
and can be scaled globally into $[0,1]$.
\end{proof}

\Cref{s3:copy-figure} shows the physical full and half gadgets and
their separate zero-port cycles.
\begin{figure}[p]
\centering
\begin{tikzpicture}[
  x=1cm,y=1cm,font=\small,
  source/.style={draw,rounded corners=2pt,align=center,inner sep=4pt},
  output/.style={draw,align=center,minimum width=1.35cm,inner sep=4pt},
  port/.style={circle,draw,fill=white,inner sep=1.5pt},
  arc/.style={-{Stealth[length=1.7mm]},semithick},
  flow/.style={fill=white,inner sep=1pt},
  heading/.style={font=\small\bfseries,anchor=west}
]
\node[heading] at (-0.8,3.0) {(a) Full occurrence, $0\le x\le2$};
\node[source] (Mf) at (1.5,2.1) {Middle input\\$S=4,\ q=\beta/2$};
\node[output] (Af) at (0,0) {$A$\\$D=4$\\$\bar q=\beta/2$};
\node[output] (Bf) at (3,0) {$B$\\$D=4$\\$\bar q=\beta/2$};
\draw[arc] (Mf) -- node[flow,left] {$4-2x$} (Af);
\draw[arc] (Mf) -- node[flow,right] {$2x$} (Bf);
\foreach \name/\xx/\target/\val/\quality/\side in {
  fA0/-0.65/Af/x/0/left,fAb/0.65/Af/x/\beta/right,
  fB0/2.35/Bf/{2-x}/0/left,fBb/3.65/Bf/{2-x}/\beta/right} {
  \node[port] (\name) at (\xx,-1.8) {};
  \node[below] at (\name) {$q=\quality$};
  \draw[arc] (\name) -- node[flow,near start,\side] {$\val$} (\target);
}

\begin{scope}[xshift=7cm]
\node[heading] at (-0.8,3.0) {(b) Half occurrence, $0\le x\le2$};
\node[source] (Mh) at (1.5,2.1) {Middle input\\$S=3,\ q=1$};
\node[output] (Ah) at (0,0) {$A$\\$D=3$\\$\bar q=1$};
\node[output] (Bh) at (3,0) {$B$\\$D=3$\\$\bar q=1$};
\draw[arc] (Mh) -- node[flow,left] {$3-3x/2$} (Ah);
\draw[arc] (Mh) -- node[flow,right] {$3x/2$} (Bh);
\foreach \name/\xx/\target/\val/\quality/\side in {
  hA0/-0.65/Ah/x/0/left,hA3/0.65/Ah/{x/2}/3/right,
  hB0/2.35/Bh/{2-x}/0/left,hB3/3.65/Bh/{1-x/2}/3/right} {
  \node[port] (\name) at (\xx,-1.8) {};
  \node[below] at (\name) {$q=\quality$};
  \draw[arc] (\name) -- node[flow,near start,\side] {$\val$} (\target);
}
\end{scope}

\node[heading] at (-0.8,-3.05) {(c) Separate closed zero-port cycles};
\foreach \kind/\yy/\signal in {F/-5.1/x_f,H/-8.0/x_h} {
  \begin{scope}[yshift=\yy cm]
  \node[anchor=west] at (-0.8,0.35) {$\kind$:};
  \node[output,minimum width=0.7cm] (A1) at (0.6,0) {$A_1$};
  \node[output,minimum width=0.7cm] (B1) at (2.5,0) {$B_1$};
  \node[output,minimum width=0.7cm] (A2) at (6.1,0) {$A_2$};
  \node[output,minimum width=0.7cm] (B2) at (8,0) {$B_2$};
  \draw[dashed,gray] (-0.05,-0.6) rectangle (3.15,0.6);
  \draw[dashed,gray] (5.45,-0.6) rectangle (8.65,0.6);
  \node at (1.55,-0.85) {occurrence 1};
  \node at (7.05,-0.85) {occurrence 2};
  \node[source] (Z1) at (4.3,0) {$Z_1$};
  \node[source] (Z2) at (4.3,1.4) {$Z_2$};
  \draw[arc] (Z1) -- node[flow,above] {$2-\signal$} (B1);
  \draw[arc] (Z1) -- node[flow,above] {$\signal$} (A2);
  \draw[arc] (Z2) -| node[flow,pos=0.25,above] {$\signal$} (A1);
  \draw[arc] (Z2) -| node[flow,pos=0.25,above] {$2-\signal$} (B2);
  \node[align=left,anchor=west] at (9.0,0.6) {Each $Z_i$:\\$S=2,\ q=0$};
  \end{scope}
}
\end{tikzpicture}
\caption{Physical copy gadgets and their zero-port links
(\cref{s3:copy-lemma,s3:half,s3:cycles}). Here $S$ and $D$ are exact
node supply and demand, $q$ is an input quality, and $\bar q$ is an
output upper bound. Open circles mark ports whose sources are assigned
when connecting gadgets; every arrow is an input-to-output physical arc.
In (a), $\beta>0$; ordinary full gadgets use $\beta=3$.
Endpoint capacities are two, except that the half gadget's quality-three
ports have capacity one; middle-arc capacities are $S$.
The displayed flows attain the quality bounds. Exact output qualities
enforce them locally; upper bounds alone enforce them after closing the
separate cycles in (c). That panel shows two occurrences per cycle and
only their zero-port arcs; middle and positive-quality arcs are omitted.
Additional occurrences are inserted in the same manner, with distinct
outputs at every occurrence. The dashed boxes group outputs and are not
physical nodes or arcs. Full and half signals are coupled by the
separate source equation $x_f+(1-x_h/2)+(1-x_h/2)=2$, using two distinct
half occurrences, as in the proof of \cref{s3:cycles}.}
\label{s3:copy-figure}
\end{figure}


The full-only degree-four construction, the full/half degree-three
construction, and the cyclic input-degree-two construction are exact
projection refinements of the same row realization. Thus none discards
the intermediate guarantees: one exact scalar quality or two upper
attributes suffice before the cyclic refinement, and one upper scalar
quality suffices afterwards. Up to this point every construction uses
positive flow contracts. Removing them is a separate mathematical step.

\subsection{An explicit exact penalty for removing contracts}\label{s3:penalty}

The next argument preserves the entire optimal value of the contracted
hard family, up to an explicit additive constant. The later constant-data
construction instead encodes threshold feasibility and needs no such
value-preserving penalty.

\begin{lemma}[A rational polyhedral error bound]\label{s3:hoffman}
Let $Q=\{v\in\R^N:Rv\le s\}$ be nonempty, $N\ge1$, and let each row
be scaled by a positive factor to have integral coefficients of absolute
value at most $K_0\ge1$. Let $R,s$ denote the resulting scaled system. In this representation,
\begin{equation}
 \operatorname{dist}_1(w,Q)\le H\max_j(R_jw-s_j)_+,
 \qquad H=N(NK_0)^{N-1}.\label{s3:hoffman-bound}
\end{equation}
\end{lemma}
\begin{proof}
Let $v$ be the Euclidean projection of $w$ onto the closed convex set
$Q$, and put $h=w-v$. If $h=0$ the result is immediate. The projection
optimality condition puts $h$ in the cone of the active row normals.
A conic representation can be chosen with linearly independent normals:
if a supported collection is dependent, move its nonnegative coefficients
along a dependence until one coefficient becomes zero and repeat. Thus
$h=B^T\lambda$, where $B$ has $r\le N$ linearly independent active
integral rows and $\lambda\ge0$. Set
$\eta=\max_j(R_jw-s_j)_+$. Since these rows are tight at $v$,
\[
 \|h\|_2^2=\lambda^TBh\le\|\lambda\|_1\eta,
 \qquad
 \|\lambda\|_1\le\sqrt N\,\|h\|_2/\sigma_{\min}(B).
\]
Since $B$ has full row rank and integral entries, $\det(BB^T)$ is
a positive integer and therefore at least one. All singular values
of $B$ are at most $\|B\|_F\le NK_0$. Their product is at least one,
so $\sigma_{\min}(B)\ge(NK_0)^{-(r-1)}\ge(NK_0)^{-(N-1)}$.
Cancel $\|h\|_2>0$ and use $\|h\|_1\le\sqrt N\|h\|_2$.
\end{proof}

This is a coarse explicit form of a polyhedral error bound in the sense
of \citet[Section~2]{s3:hoffman1952}. Its proof provides the coefficient
control needed here without computing a best error-bound constant.

\begin{theorem}[Upper flow bounds only]\label{s3:upper-only}
Pooling threshold decision is NP-complete with one pool, one upper-bound
scalar quality, no positive lower flow bounds, input out-degree at most
two, output in-degree at most three, and two pool outlets. All flow
upper bounds are at most four and all quality data can lie in $[0,1]$.
Nonnegative input production costs and output unit revenues suffice.
This proof establishes ordinary NP-completeness and an exact equality
between optimal values before and after removing contracts.
\end{theorem}
\begin{proof}
Use the Matsui-specific construction of \cref{s3:base-bypass}, followed
by \cref{s3:half,s3:cycles}, before any final quality normalization.
Every actual intake has quality $a_i>1$ and reward $b_i>0$. Let $w$
contain all copy-network bypass flows and all actual intakes $x_i$;
exclude the anchor and the two primary outlets. Include
$T=\sum_i x_i\le2$. The exact copy subsystem is a bounded rational
polytope
\[
 Q=\{w:Aw\le d,\ Ew=e\}.
\]
Here $Aw\le d$ contains all upper bounds, nonnegativity, and the
linear quality inequalities at copy outputs; $Ew=e$ lists all exact
source supplies and gadget-output demands \emph{after} the cycles and
source splitting. Its rows have zero-one coefficients and integral
right sides in $\{0,1,2,3,4\}$, and $Ew\le e$ is included in $Aw\le d$.
No primary outlet quality inequality or pool-quality parameter occurs
in this polytope. It is nonempty: the all-zero original intake signals
satisfy every homogeneous cone row, and their complementary copies and
auxiliary averages fill all exact contracts, even if the source polytope
has no positive-throughput point.

Delete the positive lower flow bounds of the physical network. Any
resulting feasible solution has $Aw\le d$ and deficits
$\Delta=e-Ew\ge0$, $\delta=\mathbf1^T\Delta$.
Apply \cref{s3:hoffman} to the rows $A,E,-E$. Clear denominators in
$A$ if needed, leaving the integral contract rows unchanged. All scaled
$A$ residuals remain nonpositive and the largest positive contract
residual is at most $\delta$. Thus some $w'\in Q$ satisfies
\begin{equation}
 \|w'-w\|_1\le H\delta,
 \qquad \|x'-x\|_1\le H\delta.\label{s3:restore-copy}
\end{equation}
The cleared coefficient bound $K_0$ and hence $H$ have polynomial
binary encoding length; row denominator products suffice to compute
such a bound. Computing a nearest point is not part of the reduction.

Restoring the copy subsystem may violate the primary outlet constraints.
To repair this, set
\[
 T=\mathbf1^Tx,\quad S=a^Tx,\quad F(x)=S(T-1)-T.
\]
For $x\ge0,T\le2$, the primary subsystem can accommodate $x$ exactly
when $F(x)\le0$. To see this for $T>0$, its quality is $S/T\ge1$;
its maximal throughput along that ray is $1+T/S$, by the two unit
outputs and anchor dilution. Equivalently $T\le1+T/S$, or $F\le0$.
For a direct feasible realization when $T\le1$, send everything to
output 2. When $T>1$, send one unit to output 2 and $T-1$ to output 1,
adding anchor $(S/T-1)(T-1)$. The total at output 1 is
$(S/T)(T-1)\le1$ and anchor flow is between zero and one.
Zero throughput is feasible separately.

Let $a_{\max}=\max_i a_i$ and $L=3a_{\max}+1$. On the convex set
$x\ge0,\mathbf1^Tx\le2$,
\[
 \frac{\partial F}{\partial x_i}=a_i(T-1)+S-1,
 \qquad \left|\frac{\partial F}{\partial x_i}\right|\le L.
\]
The original $x$ is feasible, so $F(x')_+\le L\|x'-x\|_1$.
If $F(x')\le0$, set $x''=x'$. Otherwise $T'>1,S'>1$, and set
\[
 \theta=\frac{1+T'/S'}{T'}<1,\qquad x''=\theta x'.
\]
Its throughput is exactly the feasible radial maximum and $F(x'')=0$.
Moreover,
\begin{equation}
 \|x'-x''\|_1=T'-1-T'/S'=F(x')/S'
 \le L\|x'-x\|_1\le LH\delta.\label{s3:radial-repair}
\end{equation}
Scaling preserves every homogeneous cone row and the signal box.
Rebuild the exact copy network from $x''$ and its determined averages
and complementary ports, then fill the primary outputs as above.
This produces a feasible contracted pooling solution.

Let $b_{\max}=\max_i b_i$ and $C=b_{\max}(1+L)H$. Equations
\eqref{s3:restore-copy}--\eqref{s3:radial-repair} give the uniform bound
\begin{equation}
 b^Tx\le b^Tx''+C\delta.\label{s3:profit-repair}
\end{equation}
Choose $M_0=C+1$, and add reward $M_0$ per unit at each originally
contracted source and gadget output. The new profit is
$b^Tx+M_0\mathbf1^TEw$. Put $B_0=M_0\mathbf1^Te$.
For every relaxed feasible solution,
\[
 P_{\rm new}-B_0=b^Tx-M_0\delta
 \le b^Tx''-(M_0-C)\delta\le\OPT_{\rm contract}.
\]
Conversely every contracted feasible point is relaxed-feasible with
zero deficit. Therefore
\begin{equation}
 \OPT_{\rm new}=B_0+\OPT_{\rm contract}.\label{s3:penalty-value}
\end{equation}
A positive deficit is strictly suboptimal relative to its repair.
Threshold $B_0+K$ proves NP-hardness.

The constants $a_{\max},b_{\max},K_0,H,L,C,M_0,B_0$ are explicit
rationals of polynomial bit length. Each penalty term is an ordinary
linear source or output throughput reward, so it is a legal arc
objective. No new arc or quality row is introduced and all upper flow
bounds are unchanged. The fixed-quality certificate gives membership
in NP.

For the stated economics, use instead the private conversion-port input
rewards described after \cref{s3:base-bypass}. Write that objective as
$c^Tw$, with $\|c\|_\infty\le b_{\max}$. In the relaxed instance it
need not equal $b^Tx$. Projection to $w'$ changes it by at most
$b_{\max}H\delta$; at $w'\in Q$ it equals $b^Tx'$. Rebuilding at
$x''$ changes this value by at most $b_{\max}LH\delta$. Thus the
same $C$ proves \eqref{s3:penalty-value}. Base and source rewards are
negative input production costs, and gadget-output penalties are
revenues. Add a common charge $D_0\ge M_0+b_{\max}$ at all inputs
and the same revenue at all outputs. Conservation cancels the offset
and all costs and revenues are nonnegative. Finally divide the original
nonnegative qualities by their global maximum; all quality inequalities
are preserved and data lie in $[0,1]$. The penalty may be computed in
the unscaled coordinates. Its large magnitude does not give strong
hardness or an approximation barrier.
\end{proof}

\subsection{Constant physical data and two actual feeds}\label{s3:constant-data}

Large source coefficients can instead be stored in the topology of a
linear circuit. This yields a stronger decision result whose local
physical data all belong to fixed finite sets.

\begin{lemma}[Constant-data linear circuits]\label{s3:linear-circuits}
Every rational system of linear equations and weak inequalities on
variables in $[0,2]^m$ is the exact
projection of a polynomial-size pool-free direct blending network with
one upper-bound scalar quality, input out-degree at most two, and output
in-degree at most three. Flow bounds lie in $\{0,1,2,3,4\}$, with exact
integer supply and demand contracts allowed. Qualities and upper
specifications belong to $\{0,1,3/2,3\}$. Each original coordinate
is the flow on a designated full positive port.
\end{lemma}
\begin{proof}
Use the full and half closed cycles of \cref{s3:cycles}, with ordinary
endpoint qualities zero and three and middle qualities $3/2$ and one.
Every requested port occurrence has a distinct gadget. Copies needed
by gates below are allocated first; their zero ports are then joined
into the closed cycles, so gate fan-out never increases a source degree.
The following
quality-three sources implement gates on signals in $[0,2]$:
\begin{center}
\begin{tabular}{lll}
\toprule
Equation or inequality & Requested port flows & Source supply\\
\midrule
$v=(u+w)/2$ & $v,1-u/2,1-w/2$ & exactly $2$\\
$w=u+v$ & $u,v,2-w$ & exactly $2$\\
$u=v$ & $u,2-v$ & exactly $2$\\
$u\le v$ & $u,2-v$ & at most $2$\\
$u=0$ & $u$ & exactly $0$\\
$u=1$ & $u$ & exactly $1$\\
\bottomrule
\end{tabular}
\end{center}
Each row follows directly from the source's supply constraint. Couple
full and half cycles by the exact-two source already proved. Privately
supply every unused other-endpoint port up to its capacity.

For a dyadic coefficient $c/2^L\in[0,1)$, with binary bits
$c=\sum_{k=0}^{L-1}c_k2^k$, implement
\begin{equation}
 t_0=0,\qquad t_{k+1}=(t_k+c_kx)/2\quad(0\le k<L).
 \label{s3:binary-multiplier}
\end{equation}
The bits are used from least to most significant. Induction gives
$t_L=\sum_{k=0}^{L-1}c_kx/2^{L-k}=cx/2^L$.
Every intermediate signal remains in $[0,2]$ and there are $O(L)$ gates.
Use the zero signal for $c_k=0$; coefficients zero and one may instead
use zero or equality gates.

For a rational row, clear denominators to obtain
$\sum_i c_ix_i\le d$ with integer $c_i,d$. Choose a power of two
$B>\sum_i|c_i|+|d|$. Compute each $(|c_i|/B)x_i$ and the constant
$|d|/B$ from the unit signal. Put the positive terms on one side and
negative terms on the other, placing the constant according to its sign,
and sum with addition gates. Each total and every partial sum is at
most $(2\sum|c_i|+|d|)/B<2$, so no hidden restriction is introduced
by the signal box. One comparison enforces the row; an equality gate
between the sums enforces an equality. Empty sides use zero. All these
circuits have size polynomial in the row's binary length, including
large denominators after clearing. Conversely the gate equations force
the stated row in every physical solution.

A three-port average or coupling source has capacities $(2,1,1)$ and
supply two; a three-port addition source has capacities $(2,2,2)$ and
supply two. Apply the complement-flow splitting of \cref{s3:cycles}
to each, with collector demand $\sum c_h-2$, respectively two or four.
The new input supplies and arc capacities are at most two. Collectors
have redundant upper quality three and degree three. All inputs now
have degree at most two. Full and half output demands and middle
supplies are four and three, and every other flow bound is an integer
between zero and four. Reserve a full positive port for each original
coordinate and supply it privately; this adds no constraint on its
signal. The two projection directions proved above complete the result.
\end{proof}

The lemma is a representation theorem for linear polyhedra, not a
hardness claim about LP feasibility. Dividing qualities by three gives
the fixed normalized alphabet $\{0,1/3,1/2,1\}$. It also completes the
binary-coefficient representation that the bounded-coefficient row
construction alone did not establish.

\begin{theorem}\label{s3:constant-thm}
Pooling threshold decision is strongly NP-complete with one pool, one
upper-bound scalar quality, exactly two actual feed arcs and two actual
outlet arcs, input out-degree at most two, and output in-degree at most
three. Every flow lower bound is zero and every upper bound belongs to
$\{0,1,2,3,4\}$. Qualities and upper specifications belong to
\begin{equation}
 \{0,1/66,1/33,1/22,1/11,1/2,1\}.\label{s3:palette}
\end{equation}
Input unit production costs belong to $\{0,1\}$, output unit revenues
to $\{1,2\}$, and the profit threshold is an integer at most four times
the number of inputs and outputs.
\end{theorem}
\begin{proof}
Take the source family of \cref{s3:matsui}. Write $P_4=p^{4n}$,
$s=n+n^2$, and $u_0=2P_4-p$. Let $D_0$ be the largest power of two
at most $u_0/2$, so $2D_0\le u_0<4D_0$. The generator bounds
\eqref{s3:generator-bounds}, together with $s<p^n$, imply
\[
 U_i<u_0+P_4+2P_4<5P_4,\qquad
 u_0>P_4,\qquad D_0>P_4/4.
\]
Thus $2\le U_i/D_0<20$. Also $D_0<P_4$ and
$V_i\le u_0+sp^{2n}<3P_4$, so $D_0V_i<3p^{8n}<K$.
Normalize the factors by $U\mapsto U/D_0,V\mapsto D_0V$, and define
\begin{equation}
 a_i=U_i/D_0-1\in[1,19),\quad b_i=K-D_0V_i>0,
 \quad \rho_i=(a_i-1)/32=(U_i-2D_0)/(32D_0)\in[0,1).
 \label{s3:dyadic-data}
\end{equation}
Each $b_i$ is a positive integer and each $\rho_i$ is dyadic. All
binary encodings remain polynomial; their large values will appear
only in circuit descriptions.

Let $x_i\in[0,2]$ be the source simplex signals. Use
\cref{s3:linear-circuits} to enforce the homogeneous source cone
$Cx\le0$, and introduce signals
\begin{equation}
 T=\sum_i x_i,\qquad h=\sum_i\rho_ix_i,\qquad \ell+h=T.
 \label{s3:two-feeds}
\end{equation}
All these signals lie in $[0,2]$, in particular enforcing $T\le2$.
They can be built either as the general row circuits or as direct
nonnegative partial sums and dyadic multipliers. In the latter form all
partial sums respect two whenever the intended assignment does, since
$0\le\rho_i\le1$ and $\sum_i x_i\le2$.
Only signals $\ell$ and $h$ get actual intake conversion gadgets, with
endpoint pairs $(0,1)$ and $(0,33)$ respectively. Insert each in its
signal's full cycle. Their positive endpoint \emph{qualities} permit
\cref{s3:cycles}; their flows may be zero. Their exact-two conversion
sources force actual intakes $\ell,h$. These are the only two pool feeds.
The pool's quality mass is
\begin{equation}
 \ell+33h=T+32h=\sum_i a_ix_i.\label{s3:mass-two-feeds}
\end{equation}

The pool again has capacity two and feeds two unit-capacity outlets to
unit-capacity primary outputs. The first has upper quality one and a
quality-zero anchor bypass of capacity one. The second has redundant
upper quality 33. For $T>0$ the source mixture $z=x/T$ has pool quality
$a(z)\ge1$, and its maximal radial throughput is $1+1/a(z)$.
Finally encode the affine row
\begin{equation}
 \sum_i b_ix_i\ge K\label{s3:physical-threshold}
\end{equation}
by the same binary circuits. This is implemented by sources, flows, and
quality rows, not retained as an external objective or side constraint.
It forces $T>0$ because $K>0$. A feasible contracted network therefore
gives $z\in P$ and $b(z)T\ge K$, and hence
$b(z)(1+1/a(z))\ge K$. By \eqref{s3:threshold-product} applied to
the normalized factors, this is equivalent to $U(z)V(z)\le K$ in the
original source. Conversely, a source yes witness uses
\eqref{s3:radial-max}, satisfies \eqref{s3:two-feeds}, and fills every
circuit at its prescribed signal values. Thus contracted feasibility is
exactly the source decision. Unlike the penalty construction, this
contracted network need not have a zero-intake feasible point.

Apply the three-port source splitting after all gates and both intake
conversions have been specified. It preserves both actual feeds and
outlets and gives input degree two and output degree three. The
unscaled quality alphabet, including the conversion middle qualities,
is
\[
 \{0,1/2,1,3/2,3,33/2,33\}.
\]
All upper quality specifications also lie in this set. Divide by 33 to
obtain \eqref{s3:palette}. All flow bounds are integers from zero to four.

It remains to remove positive lower flow bounds. List every positive
exact input-supply and output-demand contract of this final network as
$Ew=e$. Each $e_j\in\{1,2,3,4\}$. Retain the upper bounds and delete
all positive lower bounds. Maximize the contract-completion profit
\begin{equation}
 P_{\rm comp}=\sum_j E_jw,\qquad
 S_0=\sum_j e_j.\label{s3:completion}
\end{equation}
Every term is at most its own contract, so a feasible relaxed network
reaches $S_0$ exactly when every original contract holds. No cancellation
of deficits is possible. Consequently it reaches this threshold exactly
when the original contracted network is feasible. Each arc occurs in
at most one contracted-input term and one contracted-output term, so
arc rewards are at most two and $S_0$ is at most four times the number
of external nodes.

For the asserted economics, initially charge minus one at each contracted
input and zero at other inputs; give revenue one to each contracted
output and zero to other outputs. This realizes \eqref{s3:completion}.
Add one to every input cost and every output revenue. Conservation
cancels this offset and gives precisely the sets $\{0,1\}$ and
$\{1,2\}$. The entire transformation has polynomial size, and even
unary data encoding has polynomial size: only the threshold grows,
linearly in network size. This proves strong NP-hardness. One scalar
pool-quality parameter gives NP membership by \cref{s2:pooling-np}.
\end{proof}

Strong NP-completeness here concerns exact threshold attainment. Binary
averaging chains can distinguish extremely close flow values using
constant local coefficients. The proof gives neither an inverse-polynomial
completion gap nor an approximation barrier or robustness to a fixed
feasibility tolerance. Those conclusions do not follow from bounded
physical data alone.

\subsection{Feasibility with five contract exceptions}\label{s3:five}

Call an input exceptional if it has no exact total supply. Call an output
exceptional if it lacks either an exact demand or an exact quality.
The positive result in \cref{s5:arbitrary-exceptions} assumes bypass
degree at most two, a fixed number of exceptions, exact ordinary
contracts, and redundant common pool bounds. The next hard construction
retains the latter three conditions but permits three bypass inlets at
ordinary outputs. This is a comparison of bypass degrees: total physical
degree also counts pool arcs. The construction below has total output
degree at most three, as did the preceding construction. Positive
contracts are part of its feasibility model.

\begin{theorem}\label{s3:five-thm}
Pooling feasibility is strongly NP-complete with one pool, one scalar
quality, two actual feeds and two actual outlets, input out-degree at
most two, and output in-degree at most three. At most five external
nodes are exceptional: three inputs and two outputs. All other inputs
have exact supplies and all other outputs have exact demands and
qualities. Flow bounds lie in $\{0,1,2,3,4\}$ and quality data in
\eqref{s3:palette}. The pool has lower bound zero and redundant upper
bound two. The exceptional outputs use only upper quality specifications.
There is no economic objective or external threshold row.
\end{theorem}
\begin{proof}
Start from the contracted circuit in \cref{s3:constant-thm}, including
the circuit implementing \eqref{s3:physical-threshold}, before the
completion objective and before splitting its gate sources. Its only
nonlinear interface consists of the feeds $\ell,h$, the two primary
outputs, and the anchor. We make exact projection transformations to
remove the many auxiliary exceptions.

First replace each comparison gate $u\le v$, previously represented by
an upper-supply-two source on $u,2-v$, by an exact-supply-two source on
$u,s,2-v$, with a new signal $s\in[0,2]$. Its equation is $u+s=v$.
A feasible extension exists exactly when $u\le v$, since then
$s=v-u\in[0,2]$. Give this slack its usual copy cycle. Equality,
addition, averaging, coupling, zero, and unit gates already use exact
supplies. This includes every comparison in the source and threshold
circuits.

Second group the unused positive-quality ports by the gate that requested
their complementary ports. For an exact quality-three source requesting
flows $f_h$ with capacities $c_h$ and total supply $S$, the unused mates
carry $c_h-f_h$. Assign them to one new quality-three source of exact
supply $\sum_hc_h-S$. This introduces no further condition, since the
requested source already has $\sum_hf_h=S$. The possible supplies are
listed below:
\begin{center}
\begin{tabular}{lccc}
\toprule
Gate & Capacities & $S$ & Complement supply\\
\midrule
Average or full/half coupling & $2,1,1$ & $2$ & $2$\\
Addition or slack comparison & $2,2,2$ & $2$ & $4$\\
Equality & $2,2$ & $2$ & $2$\\
Zero & $2$ & $0$ & $2$\\
Unit & $2$ & $1$ & $1$\\
\bottomrule
\end{tabular}
\end{center}
There are no unused zero ports because their cycles are closed. There
is no dependence on a private port's supply in the proof of copying:
that proof uses only exact middle supplies, exact output demands, and
exact zero-quality links. Thus the port identities justify this grouping
before any implication from the grouped supplies is used.

Third split every three-port exact source, including the newly grouped
ones of supply four. A source of supply $S$ and capacities $c_h$ becomes
three same-quality sources of exact supplies $c_h$, each feeding its
original port $f_h$ and a new collector with $c_h-f_h$. The collector
has exact demand $\sum_hc_h-S$ and exact quality equal to that source
quality. Its demand enforces exactly $\sum_hf_h=S$. For both supply-two
and supply-four gates, collector demands are two or four. Every new
source has degree two and the collector degree three. Sources with one
or two ports already respect the input-degree bound.

All full and half cycle outputs already attain their upper quality
bounds by \cref{s3:cycles}, including conversion gadgets with positive
endpoint qualities one and 33. We may therefore require their qualities
exactly, without removing any feasible point. Collector qualities are
also automatic, because every collector receives a single source
quality. All these outputs already have exact demands.

After grouping, the only unused positive-endpoint ports that still
need variable private inputs are the positive ports in the two intake
conversion gadgets. Their complementary ports share exact-supply-two
sources with the actual pool feeds. Thus exactly two private conversion
fillers and the anchor remain variable inputs. The two primary outputs
retain variable demands and upper qualities. Every other external node
has the required exact contracts, giving five exceptions. Original source
signals need no extra designated output ports for this reduction; each
ordinary port is allocated to a gate or its full/half coupling.
The final external-node inventory is as follows. Degrees count all
physical arcs, and quality values here are before division by 33.
\begin{center}
\small
\begin{tabular}{p{0.32\linewidth}p{0.51\linewidth}c}
\toprule
External node & Final contract & Degree\\
\midrule
Middle inputs & Exact supply $4$ (full) or $3$ (half) & $2$\\
Zero-link inputs & Exact supply $2$, quality $0$ & $2$\\
Gate and complement inputs after splitting & Exact supply in
$\{0,1,2\}$, quality $3$ & $\le2$\\
Two conversion feed inputs & Exact supply $2$, quality $1$ or $33$ & $2$\\
Two private conversion fillers & Supply in $[0,2]$, quality $1$ or $33$ & $1$\\
Anchor & Supply in $[0,1]$, quality $0$ & $1$\\
Full and conversion outputs & Exact demand $4$, exact midpoint quality & $3$\\
Half outputs & Exact demand $3$, exact quality $1$ & $3$\\
Collectors & Exact demand $2$ or $4$, exact quality $3$ & $3$\\
Two primary outputs & Demands in $[0,1]$;
upper qualities $1$ and $33$ & $2,1$\\
\bottomrule
\end{tabular}
\end{center}
All input qualities are fixed. Only the two private fillers, the anchor,
and the two primary outputs lack the required exact contracts. Ordinary
outputs have three bypass inlets and no pool inlet; the primary outputs
receive the two actual pool outlets.

The transformations preserve all signal constraints and the physical
interface in both directions. In particular the threshold circuit still
excludes zero original intake, and feasibility remains equivalent to
$UV\le K$. Each change has size linear in the affected ports. Allocate a separate
gadget for every requested gate port before grouping its unused mate;
this makes the grouping disjoint and avoids parallel physical arcs. The
quality alphabet is unchanged, normalized by division by 33, and all
flow bounds remain in $\{0,1,2,3,4\}$. The two unit-capacity outlets
alone imply throughput at most two, making the pool bound redundant.
The reduction has polynomial size with fixed finite data and is strongly
NP-hard. The same fixed scalar parameter gives NP membership. No
objective or threshold row is imposed outside the physical network.
\end{proof}

\subsection{Interpretation of the restrictions}\label{s3:scope}

The two main reduction families explain different obstacles to
tractability. In the all-degree-two family, every feasible flow admits
an integral pure-mode replacement, but choosing compatible modes is
hard. Fractional mixing is therefore unnecessary for its optimal values.
In the fixed-pool family, the nonlinear interface has constant dimension,
but the surrounding linear constraints can still encode a hard
positive-product problem. This is why a fixed number of nonlinear
parameters supplies an NP certificate without necessarily supplying a
polynomial optimization algorithm.

The source of the constraints matters. Many attributes suffice on the
two-pool tree. A single upper quality suffices once bypasses carry the
linear circuit. The latter uses many external sources even though the
pool has only two actual feed arcs; it does not contradict the
no-bypass, bounded-input algorithms of
\citet{s4:boland2017,s3:haugland-hendrix2016}. The five-exception version
retains exact auxiliary supplies, demands, and qualities, and concerns
feasibility. The zero-lower-bound version always has the zero solution
and concerns attainment of an economic threshold.

Finally, the numerical conclusions differ. The represented-family FPTAS
approximates the ordinary-hard positive-product realization, whereas the
unit-data all-degree-two family has a constant approximation gap. The
finite-data bypass construction proves strong hardness of exact
threshold attainment without proving such a gap. Keeping these
conclusions separate identifies what each theorem says about exact
algorithms, certificates, and approximate optimization.
